\section*{Capítulo \ref{ch:b1:comparative-planetology} --- Planetología comparada}

\begin{solution}{exo:b1:comparative-planetology:1}
$V = \tfrac43\pi(\qty{6.9911e7}{m})^3 = \qty{1.4313e24}{m^3}$ y
$\bar\rho = \num{1898.13e24}/\num{1.4313e24} = \qty{1326}{kg.m^{-3}}$, solo 1.33 veces la
densidad del agua, aunque Júpiter tiene 318 veces la masa de la Tierra: está formado sobre todo
por los elementos más ligeros.
\end{solution}

\begin{solution}{exo:b1:comparative-planetology:2}
Mercurio (radio \qty{2440}{km}, \qty{5429}{kg.m^{-3}}) y Venus (\qty{6052}{km},
\qty{5243}{kg.m^{-3}}) son pequeños y más densos que la roca: \omterm{def:b1:comparative-planetology:planet-types}{planetas terrestres}, de roca y
hierro. Saturno (\qty{60268}{km}, \qty{687}{kg.m^{-3}}) es enorme y más ligero que el agua: un
\omterm{def:b1:comparative-planetology:planet-types}{gigante gaseoso}, formado sobre todo por hidrógeno y helio. Urano (\qty{25559}{km},
\qty{1270}{kg.m^{-3}}) es un gigante, pero mucho más pequeño que Saturno y casi el doble de
denso: un \omterm{def:b1:comparative-planetology:planet-types}{gigante de hielo}, cuyas aguas, metano y amoníaco pesan más que el hidrógeno.
\end{solution}

\begin{solution}{exo:b1:comparative-planetology:3}
$T_\mathrm{eq} = (2601.3\times0.23/(4\times\num{5.6704e-8}))^{1/4} =
(\num{2.638e9})^{1/4} = \qty{226.6}{K}$. Venus recibe 1.9 veces más luz solar que la Tierra,
pero sus nubes reflejan el \qty{77}{\percent} de ella: absorbe solo $2601.3\times0.23 = \qty{598}{W.m^{-2}}$, frente a $1361\times0.706 = \qty{961}{W.m^{-2}}$ de la Tierra.
\end{solution}

\begin{solution}{exo:b1:comparative-planetology:4}
De más uniforme a más concentrada: la Luna, Marte, la Tierra, Júpiter, Saturno. Una masa
puntual en el centro tiene $I = 0$, y por tanto $C = 0$.
\end{solution}

\begin{solution}{exo:b1:comparative-planetology:5}
$T = 15.945421\times86400 = \qty{1377684}{s}$;
$GM = 4\pi^2(\num{1.22187e9})^3/\num{1377684}^2 = \qty{3.794e16}{m^3.s^{-2}}$;
$M = \num{3.794e16}/\num{6.6743e-11} = \qty{5.685e26}{kg}$, a menos del
\qty{0.03}{\percent} de \qty{568.32e24}{kg}.
\end{solution}

\begin{solution}{exo:b1:comparative-planetology:6}
$S = 1361\times(149.6/228.0)^2 = \qty{585.9}{W.m^{-2}}$ en Marte, y
$1361\times(149.6/778.5)^2 = \qty{50.26}{W.m^{-2}}$ en Júpiter, de acuerdo con las tablas (la
pequeña diferencia en Marte procede del redondeo de las distancias medias).
\end{solution}

\begin{solution}{exo:b1:comparative-planetology:7}
$S \propto d^{-2}$ y $T_\mathrm{eq} \propto S^{1/4}$, luego $T_\mathrm{eq}
\propto d^{-1/2}$. Para Júpiter, $T_\mathrm{eq} = (50.26\times0.657/(4\times
\num{5.6704e-8}))^{1/4} = \qty{109.8}{K}$. Un planeta que radia más de lo que absorbe debe tener
una fuente interna de calor: Júpiter estaría liberando el calor almacenado desde su formación,
o producido al contraerse lentamente.
\end{solution}

\begin{solution}{exo:b1:comparative-planetology:8}
Venus: $737 - 226.6 = \qty{510}{K}$, \qty{225}{\percent} de $T_\mathrm{eq}$.
La Tierra: $T_\mathrm{eq} = (1361\times0.706/(4\sigma))^{1/4} = \qty{255.1}{K}$,
$288 - 255.1 = \qty{33}{K}$, \qty{13}{\percent}. Marte: $214 - 209.8 =
\qty{4}{K}$, \qty{2}{\percent}.
\end{solution}

\begin{solution}{exo:b1:comparative-planetology:9}
Con $\mathrm dm = 4\pi\rho r^2\,\mathrm dr$: $I = \int_0^R\tfrac23r^2\,4\pi
\rho r^2\,\mathrm dr = \tfrac{8\pi}{15}\rho R^5$, y $MR^2 = \tfrac43\pi\rho
R^5$, así que $C = \tfrac{8}{15}\cdot\tfrac34 = \tfrac25$. Una capa hueca tiene toda su masa en
$r = R$: $I = \tfrac23MR^2$, $C = \tfrac23$.
\end{solution}

\begin{solution}{exo:b1:comparative-planetology:10}
$\tfrac52C\bar\rho = 2.5\times0.35\times5429 = 4750.4$;
$x^2 = (4750.4 - 3300)/(5429 - 3300) = 1450.4/2129 = 0.6812$, $x = 0.8254$,
radio del núcleo $0.8254\times2439.7 = \qty{2014}{km}$. $x^3 = 0.5623$ y
$\rho_c = 3300 + 2129/0.5623 = \qty{7086}{kg.m^{-3}}$. El núcleo tiene
$0.5623\times7086/5429 = \qty{73}{\percent}$ de la masa. El núcleo de Mercurio
alcanza el \qty{83}{\percent} de su radio, frente a $3485/6371 =
\qty{55}{\percent}$ en la Tierra: Mercurio es sobre todo núcleo, con un manto rocoso delgado.
\end{solution}

\begin{solution}{exo:b1:comparative-planetology:11}
(a) La función $t\mapsto t^4$ es convexa, así que
$\bigl(\tfrac{T_1 + T_2}{2}\bigr)^4 \le \tfrac{T_1^4 + T_2^4}{2} =
T_\mathrm{eq}^4$, con igualdad solo si $T_1 = T_2$. (b) $T_1^4 =
2\times270^4 - 100^4 = \num{1.0529e10}$, $T_1 = \qty{320}{K}$, y la temperatura media es
$(320 + 100)/2 = \qty{210}{K}$, \qty{60}{K} por debajo de $T_\mathrm{eq}$.
\end{solution}

\begin{solution}{exo:b1:comparative-planetology:12}
$m = 4\pi R^2p/g = 4\pi\times(\num{6.0518e6})^2\times\num{9.2e6}/8.9 =
\num{4.602e14}\times\num{1.034e6} = \qty{4.76e20}{kg}$: 93 veces la masa de la atmósfera
terrestre, y \num{9.8e-5} de la masa de Venus. La Tierra guarda una cantidad comparable de
carbono, no en su aire, sino en rocas carbonatadas (calizas) y en el océano, donde el agua
líquida permitió almacenarlo.
\end{solution}

\begin{solution}{pb:b1:comparative-planetology:1}
\textbf{1.} $T = 0.31891\times86400 = \qty{27553.8}{s}$.

\textbf{2.} $GM = 4\pi^2\times(\num{9.378e6})^3/\num{27553.8}^2 =
\qty{4.289e13}{m^3.s^{-2}}$.

\textbf{3.} $M = \num{4.289e13}/\num{6.6743e-11} = \qty{6.426e23}{kg}$,
a menos del \qty{0.14}{\percent} de la \qty{6.4169e23}{kg} tabulada.

\textbf{4.} $\num{1.06e16}/\num{6.426e23} = \num{1.6e-8}$: Fobos es despreciable.

\textbf{5.} $v = 2\pi a/T = 2\pi\times\num{9.378e6}/27553.8 =
\qty{2138}{m.s^{-1}}$, unos \qty{2.1}{km.s^{-1}}.

\textbf{6.} $T = \qty{7.65}{h}$, 3.2 veces menos que la rotación de Marte: Fobos da la vuelta
más deprisa de lo que el suelo gira bajo él, así que adelanta a la rotación. Visto desde Marte,
sale por el oeste y se pone por el este.

\textbf{7.} $V = \tfrac43\pi(\num{3.3895e6})^3 = \qty{1.6312e20}{m^3}$.

\textbf{8.} $\bar\rho = \num{6.426e23}/\num{1.6312e20} = 3939$, es decir,
\qty{3940}{kg.m^{-3}}.

\textbf{9.} $3940/5514 = 0.715$. La Tierra, más grande y más masiva, comprime su interior con
más fuerza; y su núcleo de hierro es una parte mayor del planeta.

\textbf{10.} $g = GM/R^2 = \num{4.289e13}/(\num{3.3895e6})^2 =
\qty{3.73}{m.s^{-2}}$, \qty{38}{\percent} de la terrestre.

\textbf{11.} $v_\mathrm{esc} = \sqrt{2GM/R} = \sqrt{2\times\num{4.289e13}/
\num{3.3895e6}} = \qty{5.03}{km.s^{-1}}$.

\textbf{12.} $T_\mathrm{eq} = (586.2\times0.75/(4\times\num{5.6704e-8}))^{1/4}
= \qty{209.8}{K}$.

\textbf{13.} $214 - 209.8 = \qty{4}{K}$: un calentamiento por efecto invernadero débil.

\textbf{14.} $m = p/g = 636/3.733 = \qty{170}{kg.m^{-2}}$; para la Tierra
$101400/9.82 = \qty{10326}{kg.m^{-2}}$, 61 veces más.

\textbf{15.} $4\pi R^2\times170.4 = \num{1.444e14}\times170.4 =
\qty{2.46e16}{kg}$.

\textbf{16.} Sobre cada metro cuadrado hay sesenta veces menos gas que en la Tierra, así que la
absorción infrarroja, aunque sea por dióxido de carbono, es débil; y el aire de Marte contiene
muy poco vapor de agua, el otro gran absorbente.

\textbf{17.} Un Marte uniforme tendría $C = 0.4$; el valor medido, 0.366, muestra que su
densidad aumenta hacia el centro: Marte tiene un núcleo denso.

\textbf{18.} $\bar\rho = \rho_m + (\rho_c - \rho_m)x^3$.

\textbf{19.} $C = \tfrac25[\rho_m + (\rho_c - \rho_m)x^5]/\bar\rho$.

\textbf{20.} A partir de 18 y 19, $(\rho_c - \rho_m)x^3 = \bar\rho - \rho_m$ y
$(\rho_c - \rho_m)x^5 = \tfrac52C\bar\rho - \rho_m$; dividiendo se obtiene $x^2$.

\textbf{21.} $\tfrac52C\bar\rho = 2.5\times0.366\times3940 = 3605.1$;
$x^2 = (3605.1 - 3500)/(3940 - 3500) = 105.1/440 = 0.2389$, $x = 0.4887$.
$x^3 = 0.1167$ y $\rho_c = 3500 + 440/0.1167 = \qty{7269}{kg.m^{-3}}$.

\textbf{22.} $x^3\rho_c/\bar\rho = 0.1167\times7269/3940 = 0.215$, es decir, \qty{21.5}{\percent}
de la masa.

\textbf{23.} Con $\rho_m = 3450$: $x^2 = 155.1/490 = 0.3165$, $x = 0.5626$, radio del núcleo
\qty{1907}{km}, $\rho_c = \qty{6202}{kg.m^{-3}}$. Con $\rho_m =
3550$: $x^2 = 55.1/390 = 0.1413$, $x = 0.3759$, \qty{1274}{km},
$\rho_c = \qty{10894}{kg.m^{-3}}$. Un cambio del \qty{1.4}{\percent} en la densidad supuesta del
manto desplaza el radio del núcleo varios cientos de kilómetros: la densidad media y $C$ por sí
solos no pueden fijar con precisión el núcleo de Marte, y por eso se enviaron sismómetros.

\textbf{24.} $0.4887\times3389.5 = \qty{1657}{km}$: el radio del núcleo de Marte según el modelo
es \textbf{de unos \qty[text-series-to-math=true]{1660}{km}}, dentro del rango sísmico reciente
de \qty{1650}{km} a \qty{1675}{km} y muy por debajo de los \qty{1830}{km} de la tabla planetaria
--- un acuerdo que la pregunta 23 muestra que es en parte casualidad en la elección de $\rho_m$.
\end{solution}
